%%%%%%%%%%%%%%%%%%%%%%%%%%%%%%%%%%%%%%%%%%
%Dieses Template stellt die Vorgaben der IUBH an wissenschaftliche Dokumente in Latex dar
%Bitte prüfe ob alles korrekt ist, ich übernehme keine Haftung für Fehler!
%Da die IUBH auf Arial als Schriftart besteht muss xelatex zum kompilieren genutzt werden
%Erstellung 06/ 2020
%Author: Markus Rest
%Email: mgv@gmx.de
%%%%%%%%%%%%%%%%%%%%%%%%%%%%%%%%%%%%%%%%%%

%%%%%%%%%%%%%%%%%%%%%%%%%%%%%%%%%%%%%%%%%%
%%%%%%%%%%%%%%%%%%%%%%%%%%%%%%%%%%%%%%%%%%
%
% Im Idealfall muss nur in dem Unterordner  src/ editiert werden.
% Der Unterordner fixedsrc/ enthält nur Dinge, die sich selbst generieren
%
%%%%%%%%%%%%%%%%%%%%%%%%%%%%%%%%%%%%%%%%%%
%%%%%%%%%%%%%%%%%%%%%%%%%%%%%%%%%%%%%%%%%%

\documentclass[a4paper,11pt,listof=numbered,bibliography=numbered]{scrreport} % Allgemeiner Text – 11 Pkt

% Uservariablen - Art des Textes, Name, Matrikelnummer, ...
%Basics - Werte ändern
\newcommand{\IUBHStudentName}{Muster Name}
\newcommand{\IUBHStudentNumber}{12345678}
\newcommand{\IUBHStudentAddress}{Musterstraße 42, 12345 City}
\newcommand{\IUBHDateOfWork}{\today} %\today durch festes Datum mit \date{06.07.2020} ersetzen
\newcommand{\IUBHCourse}{Informatik}
\newcommand{\IUBHSupervisors}{Prof A. Bert \\ & Prof B. Chem} %Jeder weitere Betreuer muss mit "\\ & $NAME" hinzugefügt werden

%Sperrvermerk: Wenn Variable gesetzt wird die Seite gedruckt
\newcommand{\IUBHNonDisclosureCompany}{YourCompany GmbH}

%Titel der Arbeit auf Deutsch und Englisch
\newcommand{\IUBHTitleDe}{Dies ist ein guter Titel}
\newcommand{\IUBHTitleEn}{This is a good title}

%Art der Arbeit - Bitte nur ein oder auskommentieren!
%\newcommand{\IUBHTexttype}{Hausarbeit}
%\newcommand{\IUBHTexttype}{Seminararbeit}
%\newcommand{\IUBHTexttype}{Projektbericht}
%\newcommand{\IUBHTexttype}{Fallstudie}
\newcommand{\IUBHTexttype}{Bachelorarbeit}

%Benutzte Packages hier rein
\usepackage[printonlyused, withpage]{acronym} %Abkürzungsverzeichnis
\usepackage[ngerman]{babel}
\usepackage{booktabs} %Tabelle der Titelseite definieren
\usepackage{color} %Farben im Text realisieren
\usepackage{etoolbox} %IF/ELSE - Funktionalität
\usepackage{fontspec-xetex}
\usepackage[a4paper, left=2 cm, right=2 cm, top=2 cm, bottom=2 cm]{geometry}
\usepackage[automake]{glossaries}
\usepackage{graphicx} %
\usepackage{hyperref} %Links im Inhaltsverzeichnis, Abkürzungsverzeichnis usw.
\usepackage{sectsty}
\usepackage{setspace} %Abstände der Titelseite einstellen
\usepackage{svg} %Logo der IUBH einbinden

% Standards einstellen - Abstände, Schriften usw. Einstellen
% Immer:
\setmainfont{Arial} %Schrifttyp Arial  
\onehalfspacing %Zeilenabstand 1,5 

%Überschriften 12 Pkt

%Fußnoten 10 Pkt

%Variablen aus Uservorgaben setzen
\author{\IUBHStudentName}

%Glossareinträge verwalten
\makeglossaries

\newglossaryentry{latex}
{
    name=latex,
    description={Is a mark up language specially suited for scientific documents}
}

\newglossaryentry{maths}
{
    name=mathematics,
    description={Mathematics is what mathematicians do}
}

\newglossaryentry{iubh}{name={IUBH},description={IUBH Internationale Hochschule GmbH}}

\begin{document}
	\pagenumbering{Roman}
	\setcounter{tocdepth}{2}
%%%%%
%Nichtausführungen zum Text
	\begin{titlepage}
	\centering
	{\Large\textbf{\IUBHTexttype} \par}
	\vspace{2em}
	{\Huge\textbf{\IUBHTitleDe} \par}
	\vspace{0.5em}	
	{\Large\textbf{\IUBHTitleEn} \par}
	\vspace{1.5em}
	{\Large \IUBHStudentName \par}
	\vspace{1em}
	{\Large \IUBHDateOfWork \par}
	\vspace{2em}
	\includegraphics[scale=0.35]{fixedsrc/IUBH_FS_Logo_DE} \\
	\vspace{1em}
	IUBH Internationale Hochschule - Fernstudium \\
	\vspace{3em}
	\begin{tabular}{ll}
		\toprule
		Studiengang: & B.Sc. \IUBHCourse \\
		Matrikelnr.: & \IUBHStudentNumber \\
		Anschrift: & \IUBHStudentAddress \\
		Betreuer: & \IUBHSupervisors \\
		\bottomrule
	\end{tabular}	
\end{titlepage}
\newpage %Titelseite
	\expandafter\notblank\expandafter{\IUBHNonDisclosureCompany}{
	\thispagestyle{empty}
	\section*{Sperrvermerk}
	\textcolor{red}{
					Die vorliegende Arbeit beinhaltet interne und vertrauliche Informationen der Firma \IUBHNonDisclosureCompany.
					Die Weitergabe des Inhalts der Arbeit im Gesamten oder in Teilen sowie das Anfertigen von 
					Kopien oder Abschriften - auch in digitaler Form - sind grundsätzlich untersagt.\\
					Ausnahmen bedürfen der schriftlichen Genehmigung der Firma \IUBHNonDisclosureCompany.
					}	
}
{
	%Leer --> Nichts drucken!	
}
\newpage %Sperrvermerk
	\thispagestyle{empty}
\section*{Widmung}
\newpage 	%Widmung (freiwillig)
	\vspace*{\fill}
\thispagestyle{empty}
\section*{Danksagung}
An dieser Stelle möchte ich all jenen danken, die durch ihre fachliche und persönliche Unterstützung zum Gelingen dieser \IUBHTexttype\,beigetragen haben.
\vspace*{\fill}
\newpage	%Danksagung (freiwillig)
	\begin{abstract}
	\section*{Zusammenfassung}
	Das Abstract (Inhaltsangabe) besteht normalerweise aus einem Absatz, welcher die
	Hauptziele, Resultate und Schlussfolgerungen der Bachelorarbeit zusammenfasst. Es
	sollte ungefähr 200 Wörter beinhalten und nicht über eine Seite lang sein. Des Weiteren
	wird empfohlen, dass Schlüsselwörter unter den Absatz geschrieben werden.
	Schlüsselwörter sind drei bis sieben Wörter, die den Leser das Thema erkennen lassen.
	\section*{abstract}
	Dein englisches abstract hier
\end{abstract} 	%Abstract
	\tableofcontents 	%Inhaltsverzeichnis
	\newpage
	\renewcommand*{\thechapter}{\Roman{chapter}}
	\listoffigures	%Figurenverzeichnis (bei Bedarf)
	\newpage
	\listoftables %Tabellenverzeichnis
	\newpage
	\chapter{Abkürzungsverzeichnis}
%zur korrekten Formatierung die längste Abkürzung im BEGIN austauschen
\begin{acronym}[mlab]
	\acro{git}{mein Git Repository}
	\acro{mlab}{meine längste Abkürzung}
\end{acronym}
\newpage%Abkürzungsverzeichnis (bei Bedarf)
%Nichtausführungen zum Text
%%%%%

	%Vorarbeiten - Ausführungen zum Text
	\setcounter{saveromanpage}{\arabic{page}}
	\setcounter{chapter}{0}
	\setcounter{page}{3}
	\renewcommand*{\thechapter}{\arabic{chapter}}
	\pagenumbering{arabic}
   
%%%%%
%Ausführungen zum Text
    \chapter{Einleitung}
Dieser Satz kein Verb. \\
Aber dafür \ac{mlab} in meinem \ac{git} %Einleitung (normalerweise Kapitel 1 genannt)
    \chapter{Literaturanalyse}
Hier könnte deine Literaturanalyse stehen.
Außerdem testen wir hier die Glossarfunktion indem wir \gls{latex} und \gls{maths} verwenden. %Literaturanalyse (Stand der Forschung, normalerweise Kapitel 2 genannt)
    \chapter{Forschungsmethodik}

Hier soll beschrieben werden, wie denn geforscht wurde.

,,Zuweilen macht es ja wohl nichts aus, wenn man seine Arbeit auf später verschiebt.''\cite{saint-exupery__2019}%Forschungsmethodik (normalerweise Kapitel 3 genannt)
    \chapter{Forschungsergebnisse}

Hier werden die Ergebnisse dargestellt.%Forschungsergebnisse (normalerweise Kapitel 4 genannt)
    \chapter{Schlussfolgerung}

Was zeigt das ganze? Was denkst du dazu?%Schlussfolgerung (normalerweise Kapitel 5 genannt)
    \chapter{Empfehlungen}

Ein seperates "was lernen wir daraus?%Empfehlungen/Einschränkungen (kann ein separates Kapitel sein)	
% Ende Ausführungen zum Text
%%%%%
    
        %Vorarbeiten - Nichtausführungen zum Text
	\renewcommand*{\thechapter}{\Roman{chapter}}
    \setcounter{chapter}{3}
	\pagenumbering{Roman}
	\setcounter{page}{11}

    
%%%%
% Nichtausführungen zum Text
	%Quellenverzeichnis
	\bibliography{src/references}{}
	\bibliographystyle{apalike}
	%Anhang
	\chapter{Anhang}
	
	Hier können alle Anhänge eingetragen werden
	
\newpage
	%Glossar (freiwillig)
	\printglossary
	\newpage
	%Eidesstattliche Erklärung (auch Echtheitserklärung genannt)
	\chapter{Eidesstattliche Erklärung}

Hiermit versichere ich an Eides statt, dass ich die vorliegende \IUBHTexttype\ mit dem Titel ,,\IUBHTitleDe'' selbstständig und ohne Inanspruchnahme fremder Hilfe angefertigt habe. Ich habe dabei nur die angegebenen Quellen und Hilfsmittel verwendet und die aus diesen wörtlich oder inhaltlich entnommenen Stellen als solche kenntlich gemacht. Die Arbeit hat in gleicher oder ähnlicher Form noch keiner anderen Prüfungsbehörde vorgelegen. Ich erkläre mich damit einverstanden, dass die Arbeit mit Hilfe eines Plagiatserkennungsdienstes auf enthaltene Plagiate überprüft wird.

\vspace{15ex}

\IfFileExists{src/signature.png}
	{\includegraphics[width=0.3\linewidth]{src/signature}\\}
	{  }
( \IUBHStudentName\ )\\
Unterschrift
% Nichtausführungen zum Text
%%%%%
\end{document}
